\section{Method: Online RL Training}

\paragraph{Overview.}
We study online reinforcement learning for reasoning with a Qwen3-4B policy fine-tuned with LoRA. Our online RL algorithm is Group Relative Policy Optimization (GRPO), applied to two tasks: mathematical problem solving and Game of 24. Across both tasks, we keep the same high-level training loop: sample multiple rollouts per prompt, score them with a task reward, compute a group-normalized response-level advantage, and optionally add a token-level intrinsic term before the actor update.

The main goal of the study is not only to compare standard reward shaping baselines, but also to test whether token-level intrinsic signals can improve reasoning optimization. We therefore compare three classes of methods: (1) vanilla GRPO with only the task reward, (2) reward-side length baselines that modify the scalar task reward, and (3) intrinsic baselines that add a token-level signal to the GRPO advantage.

\paragraph{Base policy and optimization setup.}
For all online RL experiments, the actor policy is initialized from \texttt{Qwen3-4B} and trained with LoRA (\texttt{rank}=64, \texttt{alpha}=64). We use fp16 training and a learning rate of \(10^{-6}\). Unless otherwise noted, all runs use \texttt{rollout\_batch\_size}=16, \texttt{global\_batch\_size}=16, and \(n=8\) rollouts per prompt during training, with validation every 5 steps using one sampled rollout per prompt.

For the final Math suite, we use \texttt{max\_steps}=150 and \texttt{max\_response\_length}=1024. For the final Game of 24 suite, we use \texttt{max\_steps}=100 and \texttt{max\_response\_length}=512, since the Game of 24 training set is much smaller and responses are substantially shorter.

\paragraph{GRPO training objective.}
Let a prompt \(x_i\) produce \(n\) sampled responses \(\{y_{i,j}\}_{j=1}^n\). The reward function produces a token-level reward tensor, but under the default GRPO path used in this project, this reward is first collapsed to a scalar sequence score:
\[
R_{i,j} = \sum_t r_{i,j,t}.
\]
GRPO then normalizes these scores within the rollout group for the same prompt:
\[
\tilde{A}_{i,j} = \frac{R_{i,j} - \mu_i}{\sigma_i + \varepsilon},
\]
where \(\mu_i\) and \(\sigma_i\) are the mean and standard deviation of the \(n\) sequence scores for prompt \(x_i\). This scalar group-relative advantage is then broadcast over all valid response tokens:
\[
A^{\mathrm{ext}}_{i,j,t} = \tilde{A}_{i,j}.
\]
The actor is then updated with the standard PPO-style clipped policy gradient objective over tokens, using \(A^{\mathrm{ext}}_{i,j,t}\) or its intrinsic-augmented variant described below.

\paragraph{Math reward.}
For Math, the base reward combines answer correctness and output format. Responses are instructed to contain internal reasoning in \texttt{<think>...</think>} and a final boxed answer in \(\backslash boxed\{\}\). Let \(r_{\mathrm{acc}}\in\{0,1\}\) denote exact answer correctness, computed by extracting the boxed answer and grading it against the reference answer, and let \(r_{\mathrm{fmt}}\in\{0,1\}\) denote whether the response follows the expected output format. The external Math reward is
\[
r_{\mathrm{math}} = (1-\lambda_{\mathrm{fmt}})\, r_{\mathrm{acc}} + \lambda_{\mathrm{fmt}}\, r_{\mathrm{fmt}},
\]
with \(\lambda_{\mathrm{fmt}} = 0.1\).

\paragraph{Game of 24 reward.}
For Game of 24, each prompt contains four card values and asks the model either to produce a valid arithmetic expression that uses each card exactly once and evaluates to 24, or to output \texttt{NO} if the instance is unsolvable. As in Math, we also include a small format component encouraging \texttt{<think>...</think><answer>...</answer>} outputs. Let \(r_{\mathrm{acc}}\in\{0,1\}\) be 1 if (i) the prompt is solvable and the model outputs a valid equation using exactly the provided cards once each and evaluating to 24, or (ii) the prompt is unsolvable and the model outputs \texttt{NO}. Let \(r_{\mathrm{fmt}}\in\{0,1\}\) be the format indicator. The external Game of 24 reward is
\[
r_{24} = (1-\lambda_{\mathrm{fmt}})\, r_{\mathrm{acc}} + \lambda_{\mathrm{fmt}}\, r_{\mathrm{fmt}},
\]
again with \(\lambda_{\mathrm{fmt}} = 0.1\).

\paragraph{Reward-side length baselines.}
We include three scalar reward baselines that modify the external reward before GRPO computes the group-relative advantage.

\emph{Correct-only group-length baseline.}
This baseline is adapted from a recent paper. For each prompt group, we consider only the correct responses and compute a z-score of their response lengths. If a response is incorrect, its reward is zero and it does not participate in the z-score calculation. If a response is correct, its reward is
\[
r = 1 - \alpha \, \sigma(z_{\ell}),
\]
where \(\sigma(\cdot)\) is the sigmoid function and \(z_{\ell}\) is the per-prompt z-scored response length among correct responses only. We use \(\alpha = 0.01\).

\emph{L1-Exact.}
This baseline follows the exact-length objective:
\[
r = \mathbb{I}(y = y^\star) - \alpha \, |n_{\mathrm{target}} - n_y|,
\]
where \(n_y\) is the generated response length and \(n_{\mathrm{target}}\) is a prescribed target length. We use \(\alpha = 0.01\). For Math we set \(n_{\mathrm{target}} = 512\); for Game of 24 we set \(n_{\mathrm{target}} = 256\).

\emph{L1-Max.}
This baseline follows the soft maximum-length formulation:
\[
r = \mathbb{I}(y = y^\star)\cdot \mathrm{clip}\big(\alpha (n_{\mathrm{target}} - n_y) + \delta,\; 0,\; 1\big),
\]
where \(\delta = 0.5\). As above, we use \(\alpha = 0.01\), with \(n_{\mathrm{target}} = 512\) for Math and \(n_{\mathrm{target}} = 256\) for Game of 24.

\paragraph{Intrinsic baselines.}
Our intrinsic baselines modify the token-level advantage rather than the scalar external reward. Let \(A^{\mathrm{ext}}_{i,j,t}\) denote the broadcast GRPO advantage described above. We define a token-level intrinsic signal \(A^{\mathrm{int}}_{i,j,t}\), normalize it within each prompt group, and then combine it with the external advantage as
\[
A_{i,j,t} = A^{\mathrm{ext}}_{i,j,t} + \eta \, g_{i,j} \, A^{\mathrm{int}}_{i,j,t}.
\]
Here \(\eta\) is the intrinsic weight and \(g_{i,j}\) is an outcome-dependent gate. In our final experiments we use an asymmetric gate:
\[
g_{i,j} =
\begin{cases}
1, & \text{if the response is incorrect}, \\
-1, & \text{if the response is correct}.
\end{cases}
\]
This means that the intrinsic term is added on failures and subtracted on successes.

\paragraph{Prompt-group z-score normalization.}
The key normalization choice in the final experiments is prompt-group normalization. For each prompt, we pool all valid token-level intrinsic source values across the \(n\) rollouts for that prompt, compute a single mean and standard deviation from that pooled set, and z-score every token in the group using those shared statistics:
\[
z_{i,j,t} = \frac{s_{i,j,t} - \mu_i}{\sigma_i + \varepsilon},
\]
where \(s_{i,j,t}\) is the raw intrinsic source at token \(t\), and \(\mu_i,\sigma_i\) are computed from the pooled token values across all rollouts for prompt \(x_i\). This makes the intrinsic signal relative within a prompt group, in the same spirit as GRPO's response-level group normalization.

\paragraph{Transform choices.}
We evaluate two transforms on top of the group-zscored intrinsic signal:
\[
A^{\mathrm{int}}_{i,j,t} = z_{i,j,t}
\]
for the identity variant, and
\[
A^{\mathrm{int}}_{i,j,t} = \tanh(z_{i,j,t})
\]
for the tanh variant. The tanh transform preserves sign while smoothly limiting extreme magnitudes.

\paragraph{LeaRS: latent AR error.}
For LeaRS, we first capture a latent representation from the actor log-probability path. A lightweight autoregressive model is then trained on buffered latent trajectories to predict the next latent state. The per-token latent prediction error is used as the intrinsic source \(s_{i,j,t}\):
\[
s_{i,j,t}^{\mathrm{LeaRS}} = \| h_{i,j,t+1} - \hat{h}_{i,j,t+1} \|.
\]
We apply temporal smoothing and then prompt-group z-scoring before combining the signal with the external advantage. For the final Math suite, the LeaRS auxiliary AR model uses a lightweight configuration (\texttt{d\_model}=256, 2 layers, 4 heads) with a lighter update schedule than our early experiments (\texttt{train\_steps}=30, \texttt{train\_every\_n\_steps}=5).

\paragraph{Response-length intrinsic.}
This baseline uses the scalar response length as the intrinsic source. For each response, its length is repeated over all valid response tokens and then normalized with prompt-group z-scoring:
\[
s_{i,j,t}^{\mathrm{len}} = n_{i,j}.
\]
Although the raw source is constant within a response, its normalized and gated contribution is still meaningful as an intrinsic control signal at the advantage level.

\paragraph{Entropy intrinsic.}
This baseline uses sampled-token surprise as a proxy for token entropy. If \(\log p_{i,j,t}\) is the log-probability of the sampled token under the policy, then
\[
s_{i,j,t}^{\mathrm{ent}} = -\log p_{i,j,t}.
\]
This quantity is high when the chosen token is surprising under the current policy and low when it is confident. We again apply prompt-group z-scoring and optionally a tanh transform before adding it to the GRPO advantage.

\paragraph{Final method sets.}
Our final Math online RL suite contains: vanilla GRPO, the correct-only group-length baseline, L1-Exact, L1-Max, LeaRS + group z-score, LeaRS + group z-score + tanh, length intrinsic + group z-score, length intrinsic + group z-score + tanh, entropy intrinsic + group z-score, and entropy intrinsic + group z-score + tanh.

The shared Math settings are: \texttt{max\_steps}=150, \texttt{max\_response\_length}=1024, \texttt{rollout\_batch\_size}=16, \texttt{global\_batch\_size}=16, and \(n=8\).

Our final Game of 24 suite uses the same 10-method structure, replacing only the task reward with the Game of 24 correctness rule and using a shorter target length for the L1 baselines. The shared settings are: \texttt{max\_steps}=100, \texttt{max\_response\_length}=512, \texttt{rollout\_batch\_size}=16, \texttt{global\_batch\_size}=16, and \(n=8\). The reduction from 150 to 100 steps is motivated by dataset size: the Game of 24 training set contains 1638 examples, so a single pass with \texttt{rollout\_batch\_size}=16 yields only about 102 optimizer steps.

\paragraph{Summary.}
The central design choice in our online RL experiments is to separate reward-side and advantage-side interventions. Reward-side baselines change the scalar sequence reward before GRPO normalization, while intrinsic baselines preserve the task reward and instead add a token-level, outcome-gated term to the GRPO advantage. This allows us to compare whether improvements come from explicit scalar reward shaping, or from richer token-level guidance layered on top of the same response-level RL objective.
